\chapter{Общая цепная метрика форм Дирихле и относительные интенсивности}
\label{ch:v8-common-chain-dirichlet-rate-metric-gate}

\section{Что осталось после индексного выбора}

Предыдущий гейт однозначно зафиксировал отношение прямой и обратной
скоростей каждой пары переноса:
\begin{equation}
 \frac{\gamma_\uparrow}{\gamma_\downarrow}=e^{-2}.
\end{equation}
Но полный генератор всё ещё имеет шесть базовых интенсивностей
\begin{equation}
 \mathcal L_{\boldsymbol\kappa}
 =\kappa_L\mathcal L_L
 +\kappa_3\mathcal L_{SU(3)}
 +\kappa_2\mathcal L_{SU(2)}
 +\kappa_1\mathcal L_{U(1)}
 +\kappa_Q\mathcal L_{QLYR}
 +\kappa_X\mathcal L_{XLdR}.
 \label{eq:v8-six-base-intensities}
\end{equation}
Проверим, устраняет ли их одна метрика, построенная из цепной степени.

\section{Шестимерный конус KMS-форм}

Все шесть слагаемых по отдельности симметричны относительно выбранного
состояния $\rho_N\propto e^{-N_\partial}$ и сохраняют его. Их матрицы как
операторов на $M_{11}\oplus M_{10}$ линейно независимы: ранг семейного
спана равен шести. Поэтому для любых $\kappa_r>0$ формула
\eqref{eq:v8-six-base-intensities} остаётся KMS-симметричной и примитивной.

После удаления общего масштаба времени это оставляет пять относительных
параметров. Скан 64 независимых наборов весов в диапазоне
$10^{-4}\ldots10^4$ во всех случаях сохранил одномерную неподвижную
алгебру, но изменил щель от $2.85\cdot10^{-6}$ до $22.17$.

Таким образом, KMS-условие уже выполнило свою работу: оно связало две
ориентации одной стрелки. Оно не сравнивает интенсивности неэквивалентных
семейств шума.

\section{Даже зависящей только от степени метрика не единственна}

На шумовом пространстве оператор $\operatorname{ad}_{N_\partial}^2$ имеет
два релевантных уровня:
\begin{equation}
 \omega^2=0
 \quad\hbox{для }SU(3),SU(2),U(1),
 \qquad
 \omega^2=4
 \quad\hbox{для }L,QLYR,XLdR.
\end{equation}
Любая положительная спектральная функция задаёт метрику
\begin{equation}
 G_f=f(0)P_0+f(4)P_2.
 \label{eq:v8-degree-spectral-metric-family}
\end{equation}
После общей нормировки остаётся свободное отношение $f(4)/f(0)$.
Машинный скан $10^{-4}\leq f(4)/f(0)\leq10^4$ целиком состоит из
примитивных KMS-процессов; щель меняется от
$2.53\cdot10^{-6}$ до $0.4898$.

Несколько одинаково естественных функций дают разные результаты:
\begin{center}
\begin{tabular}{lcc}
\toprule
$f(x)$ & перенос:калибровка & щель \\
\midrule
$1$ & $1:1$ & $0.02125997$ \\
$1+x$ & $5:1$ & $0.09128433$ \\
$e^{-x}$ & $e^{-4}:1$ & $0.00045900$ \\
$(1+x)^{-1}$ & $1:5$ & $0.00469622$ \\
\bottomrule
\end{tabular}
\end{center}
Ни одна из этих функций не выделена цепной структурой.

Чистая коммутаторная форма $f(x)=x$ не является решением: она присваивает
нулевой вес всем калибровочным слагаемым. Неподвижная алгебра тогда имеет
размерность три вместо единицы. Чтобы восстановить калибровочную диффузию, нужно
добавить положительный ``пол'' $f(0)>0$, но его отношение к $f(4)$ снова
является свободным параметром.

\begin{theorem}[Конус допустимых метрик скоростей]
При уже выбранном отношении $e^{-2}$ полный процесс допускает
шестимерный положительный конус KMS-симметричных форм Дирихле, то есть пять
относительных свобод после факторизации общего времени. Ограничение
метрики спектральными функциями цепной степени сокращает конус до одного
относительного параметра, но не выбирает его.
\end{theorem}

\begin{nogocriterion}[Запрет объявления равных весов следствием общей метрики]
Выбор $f=1$ даёт простой воспроизводимый представитель, но не вывод. Столь
же ковариантны $1+x$, $e^{-x}$ и $(1+x)^{-1}$, причём они дают разные
спектральные щели. Равенство шести весов нельзя объявлять физическим без
дополнительного коротковременного корреляционного или микроскопического
условия.
\end{nogocriterion}

\begin{protocol}
\begin{enumerate}
 \item отношение вверх/вниз: \textbf{$e^{-2}$, сохранено};
 \item независимых семейств генератора: \textbf{6};
 \item относительных весов без общего времени: \textbf{5};
 \item почленная KMS-симметрия: \textbf{пройдена};
 \item положительные случайные сборки: \textbf{64 из 64 примитивны};
 \item уровни $\operatorname{ad}_N^2$: \textbf{$0$ и $4$};
 \item зависящей только от степени относительных свобод: \textbf{1};
 \item чистая коммутаторная метрика: \textbf{провал, ядро размерности 3};
 \item равновесная метрика $f=1$: \textbf{представитель, не вывод};
 \item единственная относительная метрика скоростей: \textbf{не получена};
 \item следующий тест: \textbf{коротковременный наклон полного
 корреляционного ядра}.
\end{enumerate}
\end{protocol}

Машинный сертификат сохранён в
\path{s2t_v8_common_chain_dirichlet_rate_metric_gate_results.json}.

\begin{thebibliography}{9}
\bibitem{v8-rate-wirth}
M.~Wirth,
\emph{The differential structure of generators of GNS-symmetric quantum
Markov semigroups}, arXiv:2207.09247.
\bibitem{v8-rate-cipriani}
F.~Cipriani,
\emph{The emergence of noncommutative potential theory},
arXiv:2012.14748.
\end{thebibliography}